I was able to narrow the work to a single core problem and make a clear connection that electrodynamics provides a unifying framework for my broader research direction. Below is the working title and abstract for the NeurIPS paper. I should have results to share by the weekend.

Title:
QED-MPC: Electrodynamic Transport for Geometric Composition of Hard Constraints in Offline Goal-Conditioned Planning

Abstract:
Offline goal-conditioned planning for embodied agents requires generating long-horizon behaviors that are feasible, safe, and aligned with task goals, while preserving the many valid ways a task can be completed. Existing offline reinforcement learning and generative control methods, including flow- and diffusion-based policies, are attractive because they can leverage large offline datasets, including real-world logs, without costly or unsafe online exploration. However, over long horizons, they often trade trajectory diversity for feasibility, particularly under heterogeneous hard constraints and distribution shift. We formalize this tension as the feasibility-diversity gap. We present QED-MPC, an offline model-predictive control framework that dispenses with a separate actor policy. Instead, it plans directly over a learned latent phase space using an electrodynamic transport model and an explicit representation of admissible trajectories. This formulation enables QED-MPC to enforce hard constraints during planning and to compose heterogeneous constraints without hand-tuned scalar penalties. Combined with KL-regularized trust-region optimization, the planner generates diverse candidate futures while remaining within the admissible set. We show that, in an appropriate limit, the framework connects to classical optimal control and differential-game formulations, and we establish guarantees on value error, recursive safety, and sensitivity to model mismatch under standard assumptions. Across i.i.d. and out-of-distribution evaluations on OGBench, D4RL, and OGMARL, QED-MPC reduces constraint violations, improves goal alignment, achieves competitive return, and produces a broader set of feasible long-horizon trajectories than prior offline planners, thereby narrowing the feasibility-diversity gap.

In this work I am trying to connect electrodynamics with admissible planning as strong substrate, along with actor free admissible planner unifying foundations of continuous control in novel way that admits higher diversity and improves intent coherence (especially hard constraints). The bigger picture is I am trying to connect it (especially Quantum Electrodynamics - QED, and different levels like light, electricity, fluids, etc.) as a new organizational principle for self organization for Language Grounding for Control.

Physical Phenomenon basis
https://link.springer.com/article/10.1007/BF01882786
